% !TEX root = main.tex
\Section{Evaluation}
\label{sec:evaluation}

The evaluation asks what the mechanical component adds to an agent, in completeness and in cost, and whether it changes which model is needed.
\emph{Arm against arm} (Sec.~\ref{sec:eval-arms}) holds the model fixed and varies the tactic: \textit{agent-only}, \textit{hybrid-guided}, and \textit{hybrid-flexible}, run under the same conditions.
\emph{Model against model} (Sec.~\ref{sec:eval-models}) repeats the design under a frontier model and a mid-tier model of the same family and asks whether, with \WP, the cheaper model reaches the results of the frontier model.
Both rounds share the corpus, the protocol, the scoring, and the cost accounting described first.

\SubSection{Setup}
\label{sec:eval-setup}

\Paragraph{Corpus}
The corpus is drawn from \Anon{Etna, the private Move code of the perpetual trading system Decibel}{the private Move code of a production perpetual trading system}: 23 targets from it, two from a public Move library that carry no upstream specification, and two written for the corpus.
A private codebase matters here: the model cannot have seen these functions, or any specification of them, during training, so a contract has to be inferred rather than recalled.
Every candidate is classified by the features (strata) its contract has to account for --- loops and accumulators, early returns and breaks, arithmetic overflow and underflow, casts and truncation, mutable references and in-place permutation, function values, quantifiers, composition through a callee's contract, and others, 38 in all.
A target is admitted only if it proves within seconds against a hand-written reference specification, and before the round three \emph{mutants} are authored per target --- patches to the implementation, one per obligation the contract has to pin --- each shown to be killed by the reference.
Twenty targets are selected to cover all 38 features across 15 modules, 10 of them with loops; three are deliberately guessable as controls, and five near-duplicates are held back.
With 3 arms and 4 replicates, a model round has \textbf{240 cells}.

\Paragraph{Arms and models}
The three tactics of Sec.~\ref{sec:skills} run as arms that differ only in the rendered plugin: the reference material is byte-identical, and in agent-only the \WP tool is absent rather than discouraged.
Two rounds use two models of one family through the same runtime, the Codex CLI at high effort: Sol~5.6, the frontier model, and Terra~5.6, the mid-tier model priced at about half of Sol's list.
Each session has a budget of 100--180k output tokens, one hour of wall time, and 40 seconds of solver time per verification condition.

\Paragraph{Scoring}
A cell reaches \emph{operational success} when a judge outside the agent's workspace confirms that the package compiles, the requested scope verifies, no weakening construct is present --- |pragma verify = false|, partial abort contracts, vacuous postconditions, invented preconditions --- the bytecode is unchanged, and the required contract categories are present.
The agent's own claim is never a score, and the agent receives only this acceptance feedback: the mutants are withheld until generation has ended, and a contract is never handed back for strengthening.
Operational success is necessary but weak: a vaguer contract passes it more easily.
Completeness --- postconditions as strong as the behavior --- has no deductive check in general; \WP guarantees it only given complete loop invariants and callee contracts, and warns when it cannot, a mechanism the agent-only arm lacks.
The second level is therefore \emph{mutation adequacy}: after the round, each withheld mutant is applied to the implementation and the prover rerun against the inferred specification.
A cell is \emph{complete} when it killed every mutant and \emph{failed} when one survived it; a cell whose scoring stayed inconclusive after one infrastructure retry is \emph{unmeasured} and counts in neither.
Every one of the 480 cells reached operational success, so the outcome of a cell is decided by the withheld mutants alone.

\Paragraph{Cost}
Costs are in US dollars at each model's published list prices, recomputed from the recorded token counters of every model request rather than taken from the harness's estimate; startup requests that do no task work are excluded.
Per million tokens, Sol~5.6 is \$4 uncached input, \$0.40 cached input, and \$20 output; Terra~5.6 is \$2, \$0.20, and \$12.
Thinking counts as output for both.
Means are task-weighted, and the 95\% intervals are percentile bootstraps that hold the 20 tasks fixed and resample the four replicate blocks within each task with the three arms kept together (100,000 draws); they describe this corpus, not a population of tasks.

\SubSection{Arm against Arm}
\label{sec:eval-arms}

Table~\ref{tab:summary} summarizes both rounds, Fig.~\ref{fig:cells} shows every cell, and Tables~\ref{tab:sol} and~\ref{tab:terra} in the appendix give every target.
Under both models the hybrids complete more cells than agent-only: 77 and 77 against 74 under Sol~5.6, and 79 and 78 against 74 under Terra~5.6.
All 16 failures occur on two targets.
|QP-part-025|, the in-place partition, is the one target whose contract has to characterize a whole data structure --- that the result is a rearrangement of the input --- and it defeats every arm under both models, hybrid-guided least under Terra~5.6 (1 cell) and most under Sol~5.6 (3 cells).
|OV-order-006| is the control that never aborts, because its |&&| short-circuits past the guards that would; agent-only stated an abort condition anyway in all four Sol~5.6 cells and two Terra~5.6 cells, and no hybrid cell missed it, since \WP derives the absence of an abort exactly.

\begin{table}[t]
\centering
\caption{Both rounds, per arm: mean cost in US dollars per cell, the hybrid/agent-only cost ratio with its 95\% bootstrap interval, cells that killed every withheld mutant, cells a mutant survived, cells whose scoring stayed inconclusive, mean input/output tokens in thousands, and mean wall time per cell in seconds.}
\label{tab:summary}
\scriptsize
\setlength{\tabcolsep}{2.5pt}
% Generated by eval_cells.py; do not edit.
\begin{tabular}{llrlrrrrr}
\toprule
model & arm & cost & vs. agent [95\% CI] & complete & failed & unmeas. & I/O k & wall s \\
\midrule
Sol~5.6 & agent-only & \$0.333 & --- & 74/80 & 4 & 2 & 250/4.9 & 119 \\
 & hybrid-guided & \$0.343 & 1.03 [0.94, 1.14] & 77/80 & 3 & 0 & 263/5.3 & 125 \\
 & hybrid-flexible & \$0.362 & 1.09 [1.00, 1.18] & 77/80 & 1 & 2 & 292/5.1 & 120 \\
\midrule
Terra~5.6 & agent-only & \$0.235 & --- & 74/80 & 5 & 1 & 372/6.4 & 124 \\
 & hybrid-guided & \$0.188 & 0.80 [0.73, 0.87] & 79/80 & 1 & 0 & 281/5.0 & 102 \\
 & hybrid-flexible & \$0.207 & 0.88 [0.78, 0.99] & 78/80 & 2 & 0 & 342/5.2 & 103 \\
\bottomrule
\end{tabular}

\end{table}

\begin{figure}[t]
\centering
% scatter/classes must sit on the plot, not in an axis style: applied through a
% style it accumulates per point and exhausts TeX's memory.
\newcommand{\cellplot}[1]{%
  \addplot[only marks, scatter, scatter src=explicit symbolic, mark size=1.25pt,
    scatter/classes={
      ok0={mark=*, ModelSol}, fail0={mark=o, ModelSol, thick}, unm0={mark=x, CellMuted, thick, mark size=1.8pt},
      ok1={mark=*, ModelTerra}, fail1={mark=o, ModelTerra, thick}, unm1={mark=x, CellMuted, thick, mark size=1.8pt}}]
    table[x index=1, y expr={\thisrowno{0} + (\thisrowno{2} - 0.5) * 0.36}, meta index=3] {#1};}
\begin{tikzpicture}
\begin{groupplot}[
  group style={group size=3 by 1, horizontal sep=2.5mm, yticklabels at=edge left},
  scale only axis, width=2.85cm, height=8.4cm,
  xmode=log, xmin=0.05, xmax=1.6, log ticks with fixed point,
  xtick={0.1,0.3,1}, xticklabels={0.1,0.3,1},
  minor xtick={0.05,0.2,0.5}, minor tick style={draw=none},
  ymin=-0.6, ymax=19.6, ytick={0,...,19}, y dir=reverse,
  ymajorgrids, grid style={CellGrid, thin},
  axis line style={CellAxis}, tick style={CellAxis},
  tick label style={font=\scriptsize},
  title style={font=\small, yshift=-1ex}]
\nextgroupplot[title=agent-only, yticklabels/.expanded={\CellTaskLabels}, yticklabel style={font=\ttfamily\scriptsize}]
\cellplot{tables/agent_only-cells.dat}
\nextgroupplot[title=hybrid-guided, xlabel={cost per cell (US\$)}, xlabel style={font=\scriptsize},
  legend style={at={(0.5,-0.11)}, anchor=north, font=\scriptsize, legend columns=-1, draw=none, /tikz/every even column/.append style={column sep=0.6em}},
  legend image post style={mark size=1.4pt}]
\cellplot{tables/hybrid_guided-cells.dat}
\addlegendimage{only marks, mark=*, ModelSol} \addlegendentry{Sol~5.6}
\addlegendimage{only marks, mark=*, ModelTerra} \addlegendentry{Terra~5.6}
\addlegendimage{only marks, mark=o, black, thick} \addlegendentry{mutant survived}
\addlegendimage{only marks, mark=x, CellMuted, thick, mark size=2pt} \addlegendentry{unmeasured}
\nextgroupplot[title=hybrid-flexible]
\cellplot{tables/hybrid_flexible-cells.dat}
\end{groupplot}
\end{tikzpicture}
\caption{Every cell of both rounds: cost in US dollars on a log scale, one panel per tactic, one row per target, and within each row the four Sol~5.6 replicates above the four Terra~5.6 replicates. A filled dot is a complete cell, a hollow one a cell a withheld mutant survived, a cross a cell whose scoring stayed inconclusive.}
\label{fig:cells}
\end{figure}

The effect on cost differs between the models.
Under Sol~5.6, hybrid-guided costs 1.03$\times$ agent-only [0.94, 1.14] and hybrid-flexible 1.09$\times$ [1.00, 1.18] in the task-weighted mean; both intervals reach parity, and the hybrids are cheaper on only 8 and 9 of the 20 targets.
The hybrids are cheaper on the small arithmetic contracts --- |VS-contrib-003|, |VS-shares-002|, |PM-curve-027| at about 0.6$\times$ --- where \WP's contract is already complete and the agent only has to simplify it, and more expensive on the targets that need invention beyond the loop invariant: |SM-select-022| at 1.75$\times$ under hybrid-guided, whose contract needs a recursive specification function over a function value, and the quantified-abort target |TL-lev-020| at 1.3--1.6$\times$.
On those targets the derived \WP output is large and is re-read on every turn.
The hybrids process a twentieth to a sixth more input per cell (263k and 292k against 250k) at similar output, and this additional input, billed at cache-read prices, accounts for the difference.

Under Terra~5.6 the hybrids are cheaper: hybrid-guided costs 0.80$\times$ agent-only [0.73, 0.87] and hybrid-flexible 0.88$\times$ [0.78, 0.99], cheaper on 14 of the 20 targets each, and both hybrids also finish faster (102 and 103 seconds against 124).
The direction of the token counts reverses: the hybrids read \emph{less} than agent-only (281k and 342k against 372k) and write less (5.0k and 5.2k against 6.4k), because the unaided model searches --- more verification calls, more repair rounds --- where the hybrids start from an exact contract.
The most expensive hybrid targets are the same three as under Sol~5.6, |SM-select-022|, |QP-part-025|, and |TL-lev-020|, at 1.2--1.3$\times$.

\SubSection{Model against Model}
\label{sec:eval-models}

Fig.~\ref{fig:cells} puts the two models side by side within each row, and the comparison is uniform: Terra~5.6 is cheaper than Sol~5.6 on 19 of the 20 targets in every arm, by 0.71$\times$ under agent-only and 0.55$\times$ and 0.57$\times$ under the hybrids in the task-weighted mean, at 231 against 228 complete cells over the round.
The exception is |BA-base-012|, the lemma target, where Terra~5.6 costs 1.2--1.5$\times$ Sol~5.6 because it needs more attempts to find the monotonicity lemma.
Since Terra's price list is about half of Sol's, the agent-only ratio of 0.71 implies that Terra spends about 40\% \emph{more} tokens per cell to reach the same result, while the hybrid ratios near 0.55 imply that with \WP it spends about the same: the tool closes the token gap between the two models.
In the figure, the replicate spread is about the same for both models in every arm, a factor of about 1.6 between the cheapest and the most expensive replicate of a typical target.
What changes with the tactic is how far the Terra replicates lie to the left of the Sol replicates: little under agent-only, and clearly on nearly every row under hybrid-guided.

Together, the two rounds show two effects of \WP.
The first is on completeness and does not depend on the model: under both, the hybrids complete more cells, never miss the control whose abort condition \WP derives exactly, and fail only on the one target whose contract needs a permutation argument that \WP cannot supply either.
The second is on cost and depends on the model.
\WP supplies exact contracts, but its output also enlarges the context the model re-reads on every turn.
Sol~5.6 finds the invariants unaided, so it pays for the additional context without saving elsewhere (1.03--1.09$\times$, parity within the intervals).
Terra~5.6 would otherwise search, and the savings outweigh the context (0.80$\times$); with \WP, it completes as many cells as Sol~5.6 at 55\% of its cost.
With four replicates and failures concentrated on two targets, the completeness differences are descriptive; the Terra~5.6 cost intervals exclude parity, the Sol~5.6 ones include it.

%%% Local Variables:
%%% mode: latex
%%% TeX-master: "main"
%%% End:
